\section{Unit Quiz --- test yourself}
{\small\color{slate}Fifteen multi-step questions. Answer before checking the key on the next page. Self-check notebook: \texttt{11\_Quizzes/Quiz\_SelfCheck.ipynb}.}\par\medskip
\begin{tcolorbox}[enhanced,breakable,colback=white,colframe=navyL,arc=2mm,boxrule=0.7pt,left=3mm,right=3mm,top=2mm,bottom=2mm]
\textbf{\color{navy}QV.01}\quad A dataset has 9 positive and 5 negative examples. What is its entropy?\hfill\btl{L3}\par\smallskip
{\small \textbf{\color{navyL}(A)}\ 0.500\\[1pt]\textbf{\color{navyL}(B)}\ 0.643\\[1pt]\textbf{\color{navyL}(C)}\ 0.940\\[1pt]\textbf{\color{navyL}(D)}\ 1.000}
\end{tcolorbox}
\begin{tcolorbox}[enhanced,breakable,colback=white,colframe=navyL,arc=2mm,boxrule=0.7pt,left=3mm,right=3mm,top=2mm,bottom=2mm]
\textbf{\color{navy}QV.02}\quad Outlook splits Play Tennis into Sunny (2+, 3--), Overcast (4+, 0--) and Rain (3+, 2--). What is Gain(S, Outlook)?\hfill\btl{L3}\par\smallskip
{\small \textbf{\color{navyL}(A)}\ 0.048\\[1pt]\textbf{\color{navyL}(B)}\ 0.151\\[1pt]\textbf{\color{navyL}(C)}\ 0.247\\[1pt]\textbf{\color{navyL}(D)}\ 0.940}
\end{tcolorbox}
\begin{tcolorbox}[enhanced,breakable,colback=white,colframe=navyL,arc=2mm,boxrule=0.7pt,left=3mm,right=3mm,top=2mm,bottom=2mm]
\textbf{\color{navy}QV.03}\quad What is the Gini impurity of a node with 9 positive and 5 negative examples?\hfill\btl{L3}\par\smallskip
{\small \textbf{\color{navyL}(A)}\ 0.230\\[1pt]\textbf{\color{navyL}(B)}\ 0.459\\[1pt]\textbf{\color{navyL}(C)}\ 0.541\\[1pt]\textbf{\color{navyL}(D)}\ 0.940}
\end{tcolorbox}
\begin{tcolorbox}[enhanced,breakable,colback=white,colframe=navyL,arc=2mm,boxrule=0.7pt,left=3mm,right=3mm,top=2mm,bottom=2mm]
\textbf{\color{navy}QV.04}\quad A perceptron has w = (w$_0$, w$_1$, w$_2$) = (--1, 0, 0) with bias input x$_0$ = 1. On x = (1, 1) with target y = 1 and $\eta$ = 1, what are the weights after one update?\hfill\btl{L3}\par\smallskip
{\small \textbf{\color{navyL}(A)}\ (--1, 0, 0)\\[1pt]\textbf{\color{navyL}(B)}\ (0, 1, 1)\\[1pt]\textbf{\color{navyL}(C)}\ (--2, --1, --1)\\[1pt]\textbf{\color{navyL}(D)}\ (1, 1, 1)}
\end{tcolorbox}
\begin{tcolorbox}[enhanced,breakable,colback=white,colframe=navyL,arc=2mm,boxrule=0.7pt,left=3mm,right=3mm,top=2mm,bottom=2mm]
\textbf{\color{navy}QV.05}\quad What is the smallest feed-forward network with threshold units that represents XOR?\hfill\btl{L4}\par\smallskip
{\small \textbf{\color{navyL}(A)}\ One perceptron\\[1pt]\textbf{\color{navyL}(B)}\ One hidden layer with 2 units\\[1pt]\textbf{\color{navyL}(C)}\ One hidden layer with 4 units\\[1pt]\textbf{\color{navyL}(D)}\ Two hidden layers with 1 unit each}
\end{tcolorbox}
\begin{tcolorbox}[enhanced,breakable,colback=white,colframe=navyL,arc=2mm,boxrule=0.7pt,left=3mm,right=3mm,top=2mm,bottom=2mm]
\textbf{\color{navy}QV.06}\quad A classifier has TP = 40, FP = 10 and FN = 20. What are its precision, recall and F1?\hfill\btl{L3}\par\smallskip
{\small \textbf{\color{navyL}(A)}\ 0.80, 0.67, 0.73\\[1pt]\textbf{\color{navyL}(B)}\ 0.67, 0.80, 0.73\\[1pt]\textbf{\color{navyL}(C)}\ 0.80, 0.80, 0.80\\[1pt]\textbf{\color{navyL}(D)}\ 0.57, 0.67, 0.62}
\end{tcolorbox}
\begin{tcolorbox}[enhanced,breakable,colback=white,colframe=navyL,arc=2mm,boxrule=0.7pt,left=3mm,right=3mm,top=2mm,bottom=2mm]
\textbf{\color{navy}QV.07}\quad A decision tree scores 99\% training accuracy and 70\% test accuracy; a depth-3 tree scores 85\% and 84\%. What is the right reading?\hfill\btl{L4}\par\smallskip
{\small \textbf{\color{navyL}(A)}\ The deep tree is better\\[1pt]\textbf{\color{navyL}(B)}\ The deep tree overfits; prefer the shallow tree or prune\\[1pt]\textbf{\color{navyL}(C)}\ Both underfit\\[1pt]\textbf{\color{navyL}(D)}\ The test set is wrong}
\end{tcolorbox}
\begin{tcolorbox}[enhanced,breakable,colback=white,colframe=navyL,arc=2mm,boxrule=0.7pt,left=3mm,right=3mm,top=2mm,bottom=2mm]
\textbf{\color{navy}QV.08}\quad Two independent rules support the same conclusion with CF$_1$ = 0.6 and CF$_2$ = 0.5. What is the combined CF (MYCIN)?\hfill\btl{L3}\par\smallskip
{\small \textbf{\color{navyL}(A)}\ 0.30\\[1pt]\textbf{\color{navyL}(B)}\ 0.55\\[1pt]\textbf{\color{navyL}(C)}\ 0.80\\[1pt]\textbf{\color{navyL}(D)}\ 1.10}
\end{tcolorbox}
\begin{tcolorbox}[enhanced,breakable,colback=white,colframe=navyL,arc=2mm,boxrule=0.7pt,left=3mm,right=3mm,top=2mm,bottom=2mm]
\textbf{\color{navy}QV.09}\quad Rule: IF fever (CF 0.9) AND cough (CF 0.7) THEN flu with rule CF 0.8. What is CF(flu)?\hfill\btl{L3}\par\smallskip
{\small \textbf{\color{navyL}(A)}\ 0.50\\[1pt]\textbf{\color{navyL}(B)}\ 0.56\\[1pt]\textbf{\color{navyL}(C)}\ 0.63\\[1pt]\textbf{\color{navyL}(D)}\ 0.72}
\end{tcolorbox}
\begin{tcolorbox}[enhanced,breakable,colback=white,colframe=navyL,arc=2mm,boxrule=0.7pt,left=3mm,right=3mm,top=2mm,bottom=2mm]
\textbf{\color{navy}QV.10}\quad Which pairing of classic expert system and chaining strategy is correct?\hfill\btl{L4}\par\smallskip
{\small \textbf{\color{navyL}(A)}\ MYCIN: forward; XCON: backward\\[1pt]\textbf{\color{navyL}(B)}\ MYCIN: backward; XCON: forward\\[1pt]\textbf{\color{navyL}(C)}\ Both backward\\[1pt]\textbf{\color{navyL}(D)}\ Both forward}
\end{tcolorbox}
\begin{tcolorbox}[enhanced,breakable,colback=white,colframe=navyL,arc=2mm,boxrule=0.7pt,left=3mm,right=3mm,top=2mm,bottom=2mm]
\textbf{\color{navy}QV.11}\quad In a dataset with 2\% positives, a model predicts 'negative' for everyone. What is its accuracy, and what should you report instead?\hfill\btl{L5}\par\smallskip
{\small \textbf{\color{navyL}(A)}\ 2\%; report accuracy\\[1pt]\textbf{\color{navyL}(B)}\ 98\%; report precision, recall, F1 and PR-AUC\\[1pt]\textbf{\color{navyL}(C)}\ 50\%; report ROC-AUC only\\[1pt]\textbf{\color{navyL}(D)}\ 100\%; nothing else needed}
\end{tcolorbox}
\begin{tcolorbox}[enhanced,breakable,colback=white,colframe=navyL,arc=2mm,boxrule=0.7pt,left=3mm,right=3mm,top=2mm,bottom=2mm]
\textbf{\color{navy}QV.12}\quad What are sigmoid(0), tanh(0) and ReLU(--2)?\hfill\btl{L3}\par\smallskip
{\small \textbf{\color{navyL}(A)}\ 0.5, 0, 0\\[1pt]\textbf{\color{navyL}(B)}\ 0, 0, --2\\[1pt]\textbf{\color{navyL}(C)}\ 1, 1, 0\\[1pt]\textbf{\color{navyL}(D)}\ 0.5, 0.5, 2}
\end{tcolorbox}
\begin{tcolorbox}[enhanced,breakable,colback=white,colframe=navyL,arc=2mm,boxrule=0.7pt,left=3mm,right=3mm,top=2mm,bottom=2mm]
\textbf{\color{navy}QV.13}\quad A loan model never sees religion or caste but splits first on pin code, and approval rates differ sharply by community. What is this called?\hfill\btl{L5}\par\smallskip
{\small \textbf{\color{navyL}(A)}\ Overfitting\\[1pt]\textbf{\color{navyL}(B)}\ Proxy discrimination\\[1pt]\textbf{\color{navyL}(C)}\ Label noise\\[1pt]\textbf{\color{navyL}(D)}\ Concept drift}
\end{tcolorbox}
\begin{tcolorbox}[enhanced,breakable,colback=white,colframe=navyL,arc=2mm,boxrule=0.7pt,left=3mm,right=3mm,top=2mm,bottom=2mm]
\textbf{\color{navy}QV.14}\quad Which expert-system component answers the user's 'WHY are you asking this?' and 'HOW did you conclude that?'\hfill\btl{L2}\par\smallskip
{\small \textbf{\color{navyL}(A)}\ Knowledge acquisition\\[1pt]\textbf{\color{navyL}(B)}\ Inference engine\\[1pt]\textbf{\color{navyL}(C)}\ Explanation facility\\[1pt]\textbf{\color{navyL}(D)}\ Working memory}
\end{tcolorbox}
\begin{tcolorbox}[enhanced,breakable,colback=white,colframe=navyL,arc=2mm,boxrule=0.7pt,left=3mm,right=3mm,top=2mm,bottom=2mm]
\textbf{\color{navy}QV.15}\quad How many leaves does the ID3 tree learned from the 14-example Play Tennis data have?\hfill\btl{L3}\par\smallskip
{\small \textbf{\color{navyL}(A)}\ 3\\[1pt]\textbf{\color{navyL}(B)}\ 4\\[1pt]\textbf{\color{navyL}(C)}\ 5\\[1pt]\textbf{\color{navyL}(D)}\ 7}
\end{tcolorbox}
\newpage
\subsection*{Answer key --- Unit V quiz}
\begin{longtable}{>{\leavevmode\bfseries\color{navy}}p{1.3cm} c p{13.2cm}}
\rowcolor{navy}\textcolor{gold}{\textbf{Q.}}&\textcolor{gold}{\textbf{Ans}}&\textcolor{gold}{\textbf{Worked explanation}}\\\endhead
QV.01&\textbf{\color{leaf}C}&--(9/14) log$_2$(9/14) -- (5/14) log$_2$(5/14) $\approx$ 0.410 + 0.530 = 0.940.\\
\rowcolor{ivory}QV.02&\textbf{\color{leaf}C}&0.940 -- (5/14)(0.971) -- (4/14)(0) -- (5/14)(0.971) $\approx$ 0.247.\\
QV.03&\textbf{\color{leaf}B}&1 -- (9/14)$^2$ -- (5/14)$^2$ = 1 -- 0.413 -- 0.128 $\approx$ 0.459.\\
\rowcolor{ivory}QV.04&\textbf{\color{leaf}B}&w$\cdot$x = --1 < 0, so $\hat y$ = 0 and the error is +1. Update w $\leftarrow$ w + (1)(1, 1, 1) = (0, 1, 1).\\
QV.05&\textbf{\color{leaf}B}&XOR is not linearly separable. Two hidden units (e.g. OR and NAND) combined by an AND output unit compute XOR.\\
\rowcolor{ivory}QV.06&\textbf{\color{leaf}A}&Precision 40/50 = 0.80; recall 40/60 $\approx$ 0.667; F1 = 2PR/(P+R) $\approx$ 0.727.\\
QV.07&\textbf{\color{leaf}B}&A large train--test gap signals high variance; the shallow tree generalises better.\\
\rowcolor{ivory}QV.08&\textbf{\color{leaf}C}&CF = CF$_1$ + CF$_2$(1 -- CF$_1$) = 0.6 + 0.5 $\times$ 0.4 = 0.80.\\
QV.09&\textbf{\color{leaf}B}&AND takes the minimum premise CF (0.7), multiplied by the rule CF: 0.7 $\times$ 0.8 = 0.56.\\
\rowcolor{ivory}QV.10&\textbf{\color{leaf}B}&MYCIN diagnoses goal-first (backward); XCON configures systems data-first (forward).\\
QV.11&\textbf{\color{leaf}B}&Accuracy is 98\% but recall is 0: the accuracy paradox. Imbalanced problems need class-sensitive metrics.\\
\rowcolor{ivory}QV.12&\textbf{\color{leaf}A}&$\sigma$(0) = 1/(1 + e$^0$) = 0.5; tanh(0) = 0; ReLU(--2) = max(0, --2) = 0.\\
QV.13&\textbf{\color{leaf}B}&Pin code correlates with protected attributes, so the model discriminates through a proxy. Audit group metrics and remove or constrain proxies.\\
\rowcolor{ivory}QV.14&\textbf{\color{leaf}C}&The explanation facility traces the rules fired (HOW) and the goal being pursued (WHY).\\
QV.15&\textbf{\color{leaf}C}&Outlook $\to$ Sunny: Humidity (High $\to$ No, Normal $\to$ Yes); Overcast $\to$ Yes; Rain: Wind (Strong $\to$ No, Weak $\to$ Yes). Five leaves.\\
\end{longtable}